\chapter{Coprimality and \texorpdfstring{$6/\pi^2$}{6/pi\^{}2}}\label{ch:coprime}

Every level of a rhythm is a coprime pair: by \cref{thm:nested}(ii) a block of $q_n$ beats holds $p_n$ hits with $\gcd(p_n, q_n) = 1$. This chapter is about coprimality itself: how common it is, how it appears as statistical independence, and where $\pi$ enters through $6/\pi^2$.

\section{The orchard}

Stand at the origin of the integer lattice and look out over an orchard with a tree at every point $(a,b)$ with $a, b \ge 1$. The tree at $(a,b)$ is visible exactly when no other tree stands on the segment from the origin to it, that is, when $\gcd(a,b) = 1$. \Cref{fig:orchard} shows the visible trees in a $40 \times 40$ corner, \OrchardForty{} of the grid. For a $1000 \times 1000$ corner the share is \OrchardThousand.

\begin{figure}[t]
\centering
\includegraphics[width=0.5\textwidth]{orchard}
\caption[The orchard of visible lattice points]{The orchard: lattice points $(a,b)$ with $1 \le a, b \le 40$. Filled points are visible from the origin ($\gcd(a,b) = 1$); open circles are hidden. The gold point $(7,1)$ is the first level of $\pi$'s rhythm, read as (beats, hits).}
\label{fig:orchard}
\end{figure}

\begin{theorem}[{\citealp[Thm.~332]{HardyWright2008}}]\label{thm:coprime}
The proportion of pairs $(a,b)$ with $1 \le a, b \le N$ and $\gcd(a,b) = 1$ tends to $6/\pi^2 = \SixPiSq\ldots$ as $N \to \infty$.
\end{theorem}

The heuristic behind it is short. A prime $p$ divides both $a$ and $b$ with probability $1/p^2$, and divisibility by different primes behaves independently, so the probability that no prime divides both is
\begin{equation}\label{eq:euler-product}
  \prod_{p} \left(1 - \frac{1}{p^2}\right) = \frac{1}{\zeta(2)}, \qquad \zeta(2) = \sum_{n \ge 1}\frac{1}{n^2} = \frac{\pi^2}{6}.
\end{equation}
The first equality is Euler's product formula \citep{Euler1744}: expanding each factor $(1 - p^{-2})^{-1} = 1 + p^{-2} + p^{-4} + \cdots$ and multiplying out, unique factorization produces each $1/n^2$ exactly once. The second is Euler's solution of the Basel problem \citep{Euler1740}. The theorem makes the heuristic rigorous by counting with the Möbius function. Numerically, the product over the primes below one million is $\EulerProdMillion$, against $6/\pi^2 = \SixPiSqEight$.

So $\pi$, a number about circles, measures how often two whole numbers share no prime. The convergents of $\pi - 3$, read as pairs (beats, hits) such as $(7,1)$, $(106, 15)$ and $(113, 16)$, are all visible trees.

\section{Coprime means uncorrelated}\label{sec:corr}

The orchard looks like a matrix, and it is one. Treat ``a hit every $a$th step'' as a signal $x^{(a)}_t = 1$ if $a \mid t$ and $0$ otherwise, and compare two such signals over one common period.

\begin{proposition}[the correlation of two pulse trains]\label{prop:corr}
For integers $a, b \ge 2$, the Pearson correlation of $x^{(a)}$ and $x^{(b)}$ over a full period of $\lcm(a,b)$ steps is
\begin{equation}\label{eq:corr}
  \corr(a, b) = \frac{\gcd(a,b) - 1}{\sqrt{(a-1)(b-1)}}.
\end{equation}
In particular the correlation is zero exactly when $a$ and $b$ are coprime, and it is 1 when $a = b$.
\end{proposition}

\begin{proof}
Over a period, $x^{(a)}$ has mean $1/a$ and, being 0 or 1, variance $\frac1a(1 - \frac1a) = (a-1)/a^2$. The product $x^{(a)}_t x^{(b)}_t$ is 1 exactly when $t$ is a multiple of $\lcm(a,b) = ab/\gcd(a,b)$, so its mean is $\gcd(a,b)/(ab)$. The covariance is therefore
\[
  \frac{\gcd(a,b)}{ab} - \frac{1}{a}\cdot\frac1b = \frac{\gcd(a,b) - 1}{ab},
\]
and dividing by $\sqrt{(a-1)(b-1)}/(ab)$, the product of the standard deviations, gives \eqref{eq:corr}.
\end{proof}

Our script checks \eqref{eq:corr} against a direct computation for every pair $2 \le a \ne b \le 30$. For example $\corr(4,6) = 1/\sqrt{15} = \CorrFourSix$, which is why four against six ``keeps colliding,'' while three against four, being coprime, is uncorrelated. In the $39 \times 39$ matrix of pulse trains $2, \dots, 40$ (\cref{fig:corr}), the share of off-diagonal zeros is $\CorrZeroShare$, close to $6/\pi^2$ for the same reason as the orchard. A prime's row is zero except at its multiples: primes are the ``independent'' pulse trains of arithmetic.

\begin{figure}[t]
\centering
\includegraphics[width=0.5\textwidth]{corr_matrix}
\caption[The pulse-train correlation matrix]{The correlation matrix \eqref{eq:corr} of pulse trains with periods $2$ to $40$. White cells are zero, exactly the coprime pairs; the pattern is the orchard of \cref{fig:orchard} seen as a matrix.}
\label{fig:corr}
\end{figure}

For periods $2$ to $40$ the matrix is positive definite: its eigenvalues run from $\CorrEigMin$ to $\CorrEigMax$. Random-matrix theory is used in finance to separate real correlations among stock returns from noise \citep{Laloux1999,Plerou1999}; how the spectrum of this arithmetic correlation matrix compares with such random benchmarks is a question we leave open (Chapter~\ref{ch:discussion}). [Update: see Sections~\ref{sec:oq5} and~\ref{sec:ob5}, and Sections~\ref{sec:th-top}--\ref{sec:th-bulk}, where the extreme eigenvalues are found to leading order and the spectrum is shown to have a limiting law. \addedon{\revC}] Formula \eqref{eq:corr} itself is an elementary exercise; we did not find it stated in this form in the literature we searched, but we expect it is known, and we claim no novelty for it (Chapter~\ref{ch:new}).

\section{Hydrogen's \texorpdfstring{$1/n^2$}{1/n\^{}2}: a shared formula, not shared physics}\label{sec:hydrogen}

Hydrogen's energy levels are proportional to $-1/n^2$, and its visible spectral lines, the Balmer series, come from the differences $1/4 - 1/n^2$. In wavelengths,
\begin{equation}\label{eq:balmer}
  \lambda_n = B\,\frac{n^2}{n^2 - 4}, \qquad B = 364.56\ \text{nm},
\end{equation}
which is the formula Johann Balmer found in 1885 \citep{Balmer1885}. It gives $\lambda_3 = \BalmerAlpha$\,nm (the red H-alpha line) and $\lambda_4 = \BalmerBeta$\,nm.

Add up the reciprocal squares that set the levels, $1 + \frac14 + \frac19 + \cdots$, and the answer is $\zeta(2) = \pi^2/6$, the reciprocal of the orchard's $6/\pi^2$, and also, by \eqref{eq:euler-product}, a product over the primes (\cref{fig:zeta-two}). It is tempting to say that ``hydrogen knows about the primes.'' It does not. The match is a match of formulas: the sum of $1/n^2$ is the same series wherever it appears. No physical quantity of hydrogen is the sum $\sum 1/n^2$, and nothing in the atom depends on prime factorization. We include it because it is a good illustration of how one series can sit inside several subjects, which is the theme of Chapter~\ref{ch:connections}.

\begin{figure}[t]
\centering
\includegraphics[width=0.9\textwidth]{zeta_two}
\caption[Two roads to $\pi^2/6$]{Two roads to $\pi^2/6$: the partial sums $\sum_{n \le N} 1/n^2$ (green, the reciprocal squares behind hydrogen's levels) and the partial Euler products $\prod (1 - p^{-2})^{-1}$ over the first $N$ primes (grey). The product converges much faster.}
\label{fig:zeta-two}
\end{figure}

Section~\ref{sec:18kepler} gives a neighbouring link that is shared physics. \addedon{\revD}

